\documentclass[11pt]{article}
\usepackage[margin=1in]{geometry}
\usepackage[hidelinks]{hyperref}
\setlength{\parindent}{0pt}
\setlength{\parskip}{6pt}
\setlength{\emergencystretch}{1em}
\begin{document}
\begin{center}
{\Large\bfseries Response to Reviewers\par}
\smallskip
\emph{Auditing Hidden Reasoning Charges in Commercial LLM APIs through Token-Usage Prediction}
\end{center}

Dear Reviewers,

Thank you for your feedback on our manuscript. Our point-by-point responses follow.

\section*{Response to Reviewer 1}

\begin{quote}\small
\textbf{Recommendation}

I recommend acceptance after optional minor revision. The paper addresses a timely and practically important problem: how users can audit hidden reasoning-token charges in commercial opaque LLM services when the reasoning trace itself is not visible. The proposed PALACE framework is well motivated, technically sensible, and supported by a broad set of experiments across mathematical, coding, medical, and general reasoning tasks. I found the paper interesting and potentially useful for improving transparency in LLM API billing.

\textbf{Summary}

The submission introduces PALACE, a user-side predictive auditing framework for estimating hidden reasoning-token usage from observable prompt–answer pairs. Rather than attempting to reconstruct the hidden chain of thought, PALACE learns to predict expected reasoning-token counts under a given model configuration. It combines domain-specific GRPO adaptation with a lightweight router that selects the relevant prediction module for each query. The authors evaluate the approach on four reasoning domains and compare against CoIn, MLP-based length predictors, direct LoRA fine-tuning, classification-style prediction, and merged-domain training. The results suggest that PALACE improves request-level prediction accuracy and maintains relatively stable aggregate error, making it suitable for detecting systematic token-reporting discrepancies.

\textbf{Strengths}

\begin{itemize}
\item Important problem. The paper studies an emerging transparency issue in commercial LLM services: users may be charged for hidden reasoning tokens that they cannot directly observe or verify. This is an important and underexplored problem with clear practical relevance.
\item Practical threat model. The paper acknowledges that exact verification is impossible without provider cooperation and instead formulates a realistic predictive auditing objective. The limited-cooperation assumption, where providers release a small auxiliary dataset for calibration, is reasonable and makes the proposal more deployable than methods requiring full access to internal traces.
\item Clear methodological contribution. PALACE is more than a simple length-regression baseline. The use of semantic prompt–answer modeling, GRPO-based domain adaptation, and router-based module selection is a coherent design for handling heterogeneous reasoning domains.
\item Broad empirical evaluation. The experiments cover multiple model sizes and several domains, including math, coding, medical, and general reasoning. The comparisons with LoRA, MLP, classification, merged-domain training, and CoIn help clarify where the proposed method improves over simpler alternatives.
\item Useful aggregate-auditing perspective. The paper correctly emphasizes that billing discrepancies may be hard to assess at the single-query level but can become meaningful over large workloads. Reporting both average and aggregate error is therefore appropriate and important.
\end{itemize}
\textbf{Optional Minor Revisions and Suggestions}

The paper is strong overall. I offer the following three focused suggestions to improve the precision and reproducibility of the presentation.
\end{quote}

Thank you for highlighting the practical relevance of hidden-token billing, the value of a limited-cooperation setting, and the importance of evaluating cumulative error. Your suggestions helped us make the protocol more precise.

\paragraph{1. Evaluation metrics and auditing protocol.}
\begin{quote}\small
Clarify the evaluation metrics and auditing protocol. Please provide explicit formulas, units, and normalization for AVG and AGG error, and clearly distinguish token-count prediction accuracy from inflation-detection accuracy. For the simulated inflation experiments, a concise description of the inflation rates, decision thresholds, and false-positive/false-negative criteria would make the results easier to interpret.
\end{quote}
\textbf{Response.} The revised Methods and Supplementary Section~A define AVG as mean absolute relative request-level error and AGG as the absolute difference between predicted and reference totals divided by the reference total. Both are percentages. Pass@33 applies a 33\% tolerance to the prediction error. The detection experiment holds predictions fixed, varies the reporting increase, and sets the decision threshold on clean calibration data. Flagged clean reports are false positives. Inflated reports at or below the threshold are false negatives. Supplementary Sections~B and E give the rule and results.

\paragraph{2. Auxiliary data.}
\begin{quote}\small
Describe the auxiliary-data assumption more concretely. Since deployment depends on a provider-released calibration set, please clarify its required size and contents, whether it must be specific to each model and configuration, and how representative it needs to be. A small ablation over calibration-set size would further strengthen the practical claim, if readily available.
\end{quote}
\textbf{Response.} A reference example contains the request prompt, final answer, and a trusted reasoning-token count for the audited model and configuration. For threshold calibration, we tested 50--500 clean records and observed mean FPRs of 3--6\%, with greater variability at small calibration sizes. For estimator training, a nested data-size study on matched SFT showed improvements through the full training set in all domains. Calibration examples should reflect the deployment workload, including its domain mix and distribution of reasoning lengths. Supplementary Sections~D and E describe the reference data and report both size studies.

\paragraph{3. Uncertainty and the counting protocol.}
\begin{quote}\small
Add uncertainty information for the main results. Reporting standard deviations, confidence intervals, or multiple-seed results for the central comparisons would help establish the robustness of PALACE's gains. It would also be useful to state the tokenizer/counting protocol, since reasoning-token count is the principal prediction target.
\end{quote}
\textbf{Response.} We used 2,000 question-cluster bootstrap resamples. PALACE achieved 68.6\% Pass@33 across four domains (95\% CI [66.5, 70.7]), and matched SFT achieved \(72.2\pm0.7\%\) (mean and sample standard deviation across three independent seeds). The revision uses the released reference counts. General-domain preprocessing counts tokens in the first reasoning block with the \nolinkurl{deepseek-ai/DeepSeek-V3-0324} tokenizer. The count excludes delimiters and chat-template tokens. Supplementary Sections~D and E give the counting and evaluation protocols.

\begin{quote}\small
\textbf{Overall Assessment}

This is a good paper on an important and timely topic. The central idea is compelling: even when hidden reasoning traces are unavailable, their expected token usage may still be statistically predictable from the observable prompt and answer. PALACE provides a practical and thoughtfully designed approach to this problem, and the empirical results are promising across several reasoning domains. The suggestions above are optional refinements that would improve clarity and reproducibility; they do not alter my positive assessment of the submission.
\end{quote}

\section*{Response to Reviewer 2}
\begin{quote}\small\itshape
The paper proposes PALACE, a user-side framework that predicts hidden reasoning-token counts from prompt–answer pairs using GRPO-trained domain-specific adapters and a router, for auditing commercial LLM services. The problem is important and the formulation is novel. However, several gaps must be addressed.
\end{quote}
Thank you for identifying gaps in the deployment assumptions, detection analysis, and baseline comparisons.

\paragraph{1. Distance from real world cases.}
\begin{quote}\small\itshape
The entire evaluation uses open-model reasoning, there are no experiments comparing PALACE estimates against token counts reported by an actual commercial API.

At the same time, estimating token counts of actual commercial api raises a question: What if the commercial api doesn't give a complete trace of their reasoning. (like OpenAI or Gemini)

Either provide evidence that an auditor calibrated on model A's traces can audit model B, or clarify that the auxiliary dataset must come from the exact audited model, and discuss how realistic is the settings of this paper, since if the provider offers complete and honest traces, we can already know the token number.
\end{quote}
\textbf{Response.} PALACE can operate with an API that keeps its reasoning traces private. A reference record contains the prompt, final answer, and a trusted count aligned with the audited model version, reasoning configuration, and billing definition. Providers can share these reference counts while keeping the reasoning text private. Once the estimator and threshold are set, each audit uses only the prompt, final answer, and reported usage. We use DeepSeek-R1-family traces to establish reference counts independently of the simulated reports and thereby measure detection and false alarms. Deployment on o3 or Gemini would use a reference set with trusted counts from that service. If the reference counts already contain systematic inflation, calibration can absorb that inflation, as quantified in Supplementary Section~E.

\paragraph{2. The 33\% threshold.}
\begin{quote}\small\itshape
This author claims PALACE to be able to detect \emph{subtle} misreporting, but for the 33\% threshold, inflation below \(\sim\)33\% is undetectable at the request level.
\end{quote}
\textbf{Response.} The 33\% prediction tolerance scores a count estimate against a known reference. It differs from the decision threshold, which is calibrated using clean requests. A fixed rule that flags reports above \(1.33\) times the prediction detected 47.6\% of 30\%-inflated reports but also flagged 21.8\% of clean reports. Calibrated single-request detection was modest, whereas aggregating 50 requests detected 92.9\% of 30\% shifts with 4.9\% observed FPR. The revised Methods define the two thresholds separately, and the Discussion reports the observed gap between individual and batch detection. Supplementary Section~E reports both operating points.

\subsection*{3. Baseline adequacy.}
\paragraph{3(a). CoIn as a baseline.}
\begin{quote}\small\itshape
CoIn is a cryptographic verification protocol based on hash trees and interactive auditor queries, not a model that produces token-count estimates from prompt--answer pairs. The paper does not describe how CoIn was adapted to output point estimates. Please explain this adaptation.
\end{quote}
\textbf{Response.} Table~1 and Figure~3 report CoIn verification results at fixed token-injection rates, following \href{https://aclanthology.org/2026.findings-emnlp.363/}{Sections~4.2 and 5.1 of Sun et al. (2026)}. The injection rate is \(K/N\), with \(K\) injected tokens and \(N\) original reasoning tokens. CoIn applies a verification-score threshold to reasoning-trace evidence and returns an accept or reject decision. Its CSV accuracy is the fraction of decisions matching the test labels. On a set of injected samples at a fixed rate, this equals the fraction of injected samples correctly rejected. PALACE separately estimates reasoning-token counts from prompts and final answers. Its Pass@1 measures the fraction of greedy estimates within the specified error tolerance.

\paragraph{3(b). Strength of the predictive baselines.}
\begin{quote}\small\itshape
The remaining baselines are either weak (the MLP uses only input/output lengths) or not a clean ablation (LoRA is trained on a different data format than PALACE's GRPO modules.). Please consider a stonger baseline.
\end{quote}
\textbf{Response.} We evaluated full-parameter SFT on Qwen2.5-1.5B-Instruct with matched training data and input--output format. We also tested a LightGBM baseline using length and surface features. Across the four domains, PALACE obtained 68.6\% Pass@33, compared with 54.8\% for the domain-specific LightGBM and 49.1\% for the two-length MLP. Matched SFT obtained \(72.2\pm0.7\%\) across three seeds. PALACE minus matched SFT was \(-3.9\) percentage points (95\% CI \([-5.9,-2.0]\)) for seed 0 and \(-4.1\) points (\([-6.0,-2.2]\)) for seed 1. These comparisons show that the user-side estimator benefits from access to prompt and answer content and that matched SFT is a strong training recipe for PALACE. Supplementary Section~E reports the configurations and scores.

\section*{Response to Reviewer 3}
Thank you for the detailed suggestions on calibration, direct detection, and reproducibility.

\paragraph{1. Calibration assumption.}
\begin{quote}\small\itshape
The limited-cooperation setting is reasonable and clearly stated. Since calibration data come from the provider being audited, a short discussion of how PALACE behaves if these data are unrepresentative or biased would be helpful. Even a small sensitivity experiment would add confidence.
\end{quote}
\textbf{Response.} We tested shifts in calibration length, domain mix, and reference labels. Calibrating only on short traces raised the clean-request FPR to 19--31\%. For 100-request batches, calibrating on a uniform mixture and testing on a workload with 70\% Code raised PALACE FPR to 25.9\%. Scaling each domain using its calibration totals reduced this to 3.5\%. Inflating calibration labels by 10\% raised the threshold and substantially impaired detection of shifts up to that size. Supplementary Section~E gives these results and the correction.

\paragraph{2. Evaluating detection directly.}
\begin{quote}\small\itshape
Section 2.3 is titled around identifying reporting discrepancies, but Figure 3 mainly shows accuracy across error thresholds. The paper's central message would be stronger with a simple simulation: inflate reported counts by, say, 5--50\%, then report TPR/FPR or ROC-AUC, plus how many requests are needed to detect systematic inflation at the aggregate level.
\end{quote}
\textbf{Response.} The revised Figure~3 caption distinguishes prediction Pass@1 for PALACE and LoRA from CoIn detection rates. We now report detection separately at 10\%, 20\%, 30\%, and 50\% inflation. For the domain-specific PALACE auditor, 30\% overreporting was detected in 92.9\% of 50-request batches (4.9\% FPR) and 99.7\% of 100-request batches (5.2\% FPR). The first tested batch size above 90\% detection at 10\% inflation was 500, with 96.5\% detection and 8.4\% FPR. Supplementary Section~E reports achieved false-positive rates and the workload-size sweep.

\paragraph{3. Request-level vs.\ aggregate use.}
\begin{quote}\small\itshape
Pass@1 accuracy is strong on general reasoning and more moderate (about 60\%) on math, coding and medical tasks. It may help to position PALACE primarily as an aggregate monitoring tool. Calibrated prediction intervals (e.g., conformal methods) could make request-level flags more interpretable.
\end{quote}
\textbf{Response.} The observed detection rates support using PALACE primarily for aggregate monitoring. At 30\% overreporting, calibrated request-level detection ranged from 11.3\% to 36.8\% across domains. Nominal 90\% same-domain conformal intervals covered 89.4--91.4\% of references for PALACE. Their upper-to-lower bound ratios ranged from 2.0 to 4.6, and Code and Medical coverage fell to 83--85\% under calibration shift. Supplementary Section~E reports both coverage and interval width to characterize individual-request uncertainty.

\paragraph{4. Aggregated error.}
\begin{quote}\small\itshape
It would help to define the units of the AVG and AGG errors. The statement that PALACE consistently achieves the lowest aggregated error could be refined, since MLP and LoRA obtain lower AGG in some settings in Table 1. The authors' point that low AGG alone is insufficient is well taken and could be made explicit.
\end{quote}
\textbf{Response.} Table~1 shows that MLP or LoRA achieves lower AGG in some settings. The Results explain why aggregate error can remain low when a simple baseline predicts close to the dataset mean despite large errors on individual requests. The revised Methods and Supplementary Section~A give the formulas, denominators, and percentage units for AVG and AGG.

\paragraph{5. Generalization.}
\begin{quote}\small\itshape
All traces come from DeepSeek-R1 models. Cross-model transfer experiments, or a brief discussion of applicability to services such as o3 or Gemini, would broaden the impact. An ablation that removes the token-length cues from the prompt would also clarify the contribution of semantic modeling.
\end{quote}
\textbf{Response.} The benchmark supplies one target model per domain. To measure transfer while holding the domain fixed, the evaluation would need reference counts from a second model in the same domain. We therefore evaluate matched-model calibration. Applying this protocol to o3 or Gemini uses a model-specific reference set with counts aligned to that service's billing definition. We also tested token-length cues. Removing numeric length cues during both SFT training and testing changed Pass@33 by \(+0.4,+0.4,-0.6,+1.6\) percentage points across General, Math, Code, and Medical, with every confidence interval including zero. Removing the prompt and answer text while retaining length cues at PALACE inference reduced Pass@33 in every domain. On Math it fell from 63.8\% to 19.2\%. Supplementary Section~E reports both ablations.

\paragraph{6. Baselines and metrics.}
\begin{quote}\small\itshape
A sentence explaining how CoIn's Pass@1 was computed would help readers, as CoIn is a verification protocol. Clarifying how classification accuracy is defined in Table 2 would make that comparison clearer. A note on why Pass@5 is relevant for auditing would also be useful.
\end{quote}
\textbf{Response.} As detailed in our response to Reviewer~2, Comment~3(a), the CoIn percentages are CSV verification accuracies at fixed injection rates. On injected samples, they measure the fraction correctly rejected. Table~1 now labels the CoIn row as detection rate. Pass@1 measures single-prediction relative-error accuracy. Table~2 compares five-bucket classification with count regression. The Additional Experimental Setup defines five domain-specific buckets. The released classifier takes the argmax of the class logits and reports the fraction of examples whose predicted bucket matches the reference bucket. The revised Table~2 caption states the separate scoring criteria for classification and regression. Pass@5 measures whether sampling produces at least one accurate candidate when the true count is available to score all candidates. In the revision experiments, sampling five candidates at temperature 0.8 improved oracle accuracy by 1--3 points. Pass@5 therefore measures the oracle benefit of sampling. The detection rule uses a single greedy estimate.

\paragraph{7(a). Dataset splits, router accuracy, and run variance.}
\begin{quote}\small\itshape Report dataset split sizes, router accuracy, and the variance across runs.\end{quote}
\textbf{Response.} We split each released 1,000-example evaluation pool by question hash into 500 calibration and 500 test examples per domain. The separate Math training run in Figure~4 used 90,000 examples. The revision evaluates PALACE using the released checkpoints and their evaluation pools. The reimplemented TF-IDF plus logistic regression router achieved 99.08\% classification accuracy and was within 0.2 Pass@33 points of oracle routing. Matched SFT achieved \(72.2\pm0.7\%\) across three independently trained seeds. Each released PALACE configuration supplies one checkpoint. We use question-cluster intervals to quantify test-sample variability at those fixed weights and independent seeds to quantify run-to-run variability for SFT.

\paragraph{7(b). Training epochs.}
\begin{quote}\small\itshape Reconcile \emph{3 epochs} with the one-epoch curve in Figure 4.\end{quote}
\textbf{Response.} The revised Methods state the three-epoch default and identify the Figure~4 run as an exception. GRPO uses veRL with rollout temperature 0.8, batch size 128, eight rollouts, update batch size 128, KL coefficient 0.001, and learning rate \(10^{-6}\). The original LoRA baseline uses rank 128, batch size 128, and learning rate \(10^{-5}\). Figure~4 shows a separate Qwen2.5-3B Math learning-curve run on 90,000 examples for one epoch.

\paragraph{7(c). Loss in Equation (1).}
\begin{quote}\small\itshape Define the loss \(L\) in Eq.\ (1).\end{quote}
\textbf{Response.} The revised Methods define the training objective after Equation~(1) as next-token cross-entropy on the tokenized numeric answer. Prompt positions are masked. With answer tokens \(a_1,\ldots,a_T\), the per-example loss is \(\mathcal L=-T^{-1}\sum_{t=1}^{T}\log p_\theta(a_t\mid x,a_{<t})\), where \(x\) is the input prompt--answer pair in the auditor template.

\paragraph{7(d). GRPO domain dataset.}
\begin{quote}\small\itshape Correct the GRPO description, which refers to \(D_C\) for all domains.\end{quote}
\textbf{Response.} The revised Methods state that each GRPO adapter is trained on \(\mathcal D_d\), the dataset for its selected domain \(d\). The coding adapter uses \(\mathcal D_C\).

\paragraph{7(e). Short reasoning traces.}
\begin{quote}\small\itshape Consider how the relative-error reward in Eq.\ (2) behaves for very short reasoning traces.\end{quote}
\textbf{Response.} A one-token error reduces this relative-error reward by \(1/y\) before clipping, so the same absolute error has greater influence for shorter traces. The reduction is 1 at \(y=1\) and approximately 0.0042 at \(y=238\), the shortest reference in our test set. A zero-token reference has no relative-error denominator and calls for a separate exact-zero criterion.

\begin{quote}\small
Overall, this is a promising and well-motivated study, and addressing these points would substantially strengthen its contribution.
\end{quote}

Sincerely,\\
The authors
\end{document}
